% Options for packages loaded elsewhere
\PassOptionsToPackage{unicode}{hyperref}
\PassOptionsToPackage{hyphens}{url}
\documentclass[
  english,
  10pt,
]{article}
\usepackage{xcolor}
\usepackage[a4paper,margin=18mm]{geometry}
\usepackage{amsmath,amssymb}
\setcounter{secnumdepth}{-\maxdimen} % remove section numbering
\usepackage{iftex}
\ifPDFTeX
  \usepackage[T1]{fontenc}
  \usepackage[utf8]{inputenc}
  \usepackage{textcomp} % provide euro and other symbols
\else % if luatex or xetex
  \usepackage{unicode-math} % this also loads fontspec
  \defaultfontfeatures{Scale=MatchLowercase}
  \defaultfontfeatures[\rmfamily]{Ligatures=TeX,Scale=1}
\fi
\usepackage{lmodern}
\ifPDFTeX\else
  % xetex/luatex font selection
\fi
% Use upquote if available, for straight quotes in verbatim environments
\IfFileExists{upquote.sty}{\usepackage{upquote}}{}
\IfFileExists{microtype.sty}{% use microtype if available
  \usepackage[]{microtype}
  \UseMicrotypeSet[protrusion]{basicmath} % disable protrusion for tt fonts
}{}
\makeatletter
\@ifundefined{KOMAClassName}{% if non-KOMA class
  \IfFileExists{parskip.sty}{%
    \usepackage{parskip}
  }{% else
    \setlength{\parindent}{0pt}
    \setlength{\parskip}{6pt plus 2pt minus 1pt}}
}{% if KOMA class
  \KOMAoptions{parskip=half}}
\makeatother
\usepackage{longtable,booktabs,array}
\usepackage{caption}
\captionsetup[table]{skip=6pt}
\newcounter{none} % for unnumbered tables
\usepackage{multirow}
\usepackage{calc} % for calculating minipage widths
% Correct order of tables after \paragraph or \subparagraph
\usepackage{etoolbox}
\makeatletter
\patchcmd\longtable{\par}{\if@noskipsec\mbox{}\fi\par}{}{}
\makeatother
% Allow footnotes in longtable head/foot
\IfFileExists{footnotehyper.sty}{\usepackage{footnotehyper}}{\usepackage{footnote}}
\makesavenoteenv{longtable}
\usepackage{graphicx}
\makeatletter
\newsavebox\pandoc@box
\newcommand*\pandocbounded[1]{% scales image to fit in text height/width
  \sbox\pandoc@box{#1}%
  \Gscale@div\@tempa{\textheight}{\dimexpr\ht\pandoc@box+\dp\pandoc@box\relax}%
  \Gscale@div\@tempb{\linewidth}{\wd\pandoc@box}%
  \ifdim\@tempb\p@<\@tempa\p@\let\@tempa\@tempb\fi% select the smaller of both
  \ifdim\@tempa\p@<\p@\scalebox{\@tempa}{\usebox\pandoc@box}%
  \else\usebox{\pandoc@box}%
  \fi%
}
% Set default figure placement to htbp
\def\fps@figure{htbp}
\makeatother
\usepackage{svg}
\ifLuaTeX
\usepackage[bidi=basic,shorthands=off]{babel}
\else
\usepackage[bidi=default,shorthands=off]{babel}
\fi
\ifLuaTeX
  \usepackage{selnolig} % disable illegal ligatures
\fi
\setlength{\emergencystretch}{3em} % prevent overfull lines
\providecommand{\tightlist}{%
  \setlength{\itemsep}{0pt}\setlength{\parskip}{0pt}}
\usepackage{bookmark}
\IfFileExists{xurl.sty}{\usepackage{xurl}}{} % add URL line breaks if available
\urlstyle{same}
% fallback for those not using the hyperref driver hyperxmp:
\makeatletter
\@ifundefined{xmpquote}{\newcommand{\xmpquote}[1]{#1}}{}
\makeatother
\hypersetup{
  pdftitle={AutoE2E: Open-Loop Evaluation of a Map- and Oracle-Route-Conditioned Temporal BEV Planner},
  pdflang={en},
  hidelinks,
  pdfcreator={LaTeX via pandoc}}

\title{AutoE2E: Open-Loop Evaluation of a Map- and
Oracle-Route-Conditioned Temporal BEV Planner}
\author{}
\date{}

\begin{document}
\maketitle

\section{AutoE2E: Open-Loop Evaluation of a Map- and
Oracle-Route-Conditioned Temporal BEV
Planner}\label{autoe2e-open-loop-evaluation-of-a-map--and-oracle-route-conditioned-temporal-bev-planner}

\subsection{Authors}\label{authors}

Ryota Yamada

Amazon Web Services, Inc.\\
Autoware Foundation Robotaxi Working Group (AutoE2E),
github.com/autowarefoundation/auto\_e2e, Apache-2.0

\begin{otherlanguage}{english}

\subsection{Abstract}\label{abstract}

When a policy consumes an HD map and a route at runtime, a camera-only
benchmark omits modalities the policy was designed to use. We evaluate
AutoE2E, a UniAD-inspired but narrower temporal BEV planner, under a
runtime-modality-matched but oracle-route condition: the route corridor
and destination are reconstructed retrospectively from the logged
whole-scene trajectory. A frozen ResNet-50 BEVFormer V2 T8 camera BEV is
fused with separately encoded 14-channel map and 2-channel route rasters
through a channel-wise gate and deformable cross-attention; a GRU
predicts 64 acceleration/curvature steps. The external KITScenes
protocol supplies 4.0 s of real egomotion history, left-zero-padded to
6.4 s, and valid targets through 5.0 s. On KITScenes Val v3.5 (117
scenes, 11,035 samples), deterministic replay of published controls
gives KITScenes Epoch 5 the lowest reported means at 2, 3 and 5 s (5 s
ADE/FDE: 1.9405/5.5645 m), while Epoch 7 is lowest at 1 s. Internal
validation shows that Epoch 7 increases a censored projected-progress
proxy while increasing 6.4 s error and lowering comfort; strict
route-corridor success is 0.061--0.076 and drivable-area success
0.414--0.444. The camera-only KITScenes Test track is a different scene
population and is not a map or route ablation. No statistical
significance or causal navigation-input gain is claimed. A paired
same-scene Camera / Camera+Map / Camera+Map+Route evaluation remains
required.

\end{otherlanguage}

\begin{otherlanguage}{english}

\subsection{Keywords}\label{keywords}

end-to-end autonomous driving, temporal BEV, HD-map conditioning, route
conditioning, open-loop evaluation, KITScenes, nuPlan

\end{otherlanguage}

\subsection{Introduction}\label{introduction}

End-to-end (E2E) driving research has developed largely around
camera-only open-loop evaluation on nuScenes \hyperref[ref-nuscenes]{}.
UniAD \hyperref[ref-uniad]{}and VAD \hyperref[ref-vad]{}demonstrated
planning-oriented integrated stacks in this setting, with temporal BEV
representations such as BEVFormer \hyperref[ref-bevformer]{}providing a
shared basis. Deployed vehicle policies, however, often receive more
than cameras. The deployment contract of AutoE2E supplies an HD-map
raster and a selected-route raster at every step in addition to camera
images. Evaluating this policy only on a camera-only benchmark reports
behavior under a condition that omits information the policy was
designed to consume. Conversely, reporting only the map-and-route
condition would hide behavior when navigation inputs are unavailable.
These evaluations answer different questions and should remain separate.

This paper makes a deliberately modest claim. We have not yet run an
experiment showing that maps and routes causally improve planning
accuracy. What we do present is (i) an evaluation protocol whose primary
benchmark uses the input set that matches the deployment contract
(Camera + HD Map + Route) while retaining a camera-only condition as a
separate track, (ii) the results measured under that protocol for four
training checkpoints, (iii) implementation facts verified against the
source code and the checkpoint registration metadata, (iv) the
trade-offs among short-horizon accuracy, long-horizon accuracy, route
progress, drivable-area compliance, and comfort observed in internal
validation, and (v) the specification of the same-scene controlled
experiment that a causal claim would require.

\begin{figure}
\centering
\pandocbounded{\includegraphics[keepaspectratio,width=\linewidth]{figures_pdf/fig01_architecture.pdf}}
\caption{The evaluated AutoE2E configuration
(bevformer\_v2\_t8\_split\_navigation\_v5). At each of eight times, six
512×512 camera views are processed: the current time plus seven history
times. A 1024×1024 front view is added only at the current time. A
frozen ResNet-50 + FPN and six-layer single-frame BEVFormer V2 encoder
are applied independently per time, followed by T8 convolutional fusion
across the eight BEVs; a navigation pathway that encodes the 14-channel
map and the 2-channel route separately, gated summation, deformable
navigation fusion, 6.4 s of ego-motion history, a GRU control planner,
and unicycle integration. All reported counts are parameter counts
measured by instantiating the implementation (79,906,522 total;
1,448,582 trainable). Dashed boxes denote branches that are present in
the checkpoint but disabled.}\label{fig-arch}
\end{figure}

This paper is organized around the following five research questions
(RQs). RQ1: on the deployment-aligned KITScenes Val protocol, where
camera, HD map, and route are all available, with what accuracy does the
route-conditioned policy predict the ego trajectory? RQ2: on the same
KITScenes Val scenes, how does the KITScenes fine-tuned checkpoint
compare with the nuPlan-trained checkpoint? RQ3: how does performance
change on the camera-only KITScenes Test track, which is a different
scene population? RQ4: across training epochs, what trade-offs emerge
among short-horizon accuracy, long-horizon accuracy, route progress,
drivable-area compliance, and comfort? RQ5: given the current evidence,
what can and cannot be said about the effect of the HD map and the
route?

We limit our contributions to the following four points. First, a
source-code-verified description and an exact parameter budget for a
planning-oriented model that combines separately encoded HD map and
route rasters with a frozen BEVFormer V2 T8 camera BEV through a gated
residual (\hyperref[fig-arch]{}, \hyperref[tab-params]{}). Second, an
evaluation protocol that separates the primary benchmark, aligned with
the deployment contract, from the camera-only secondary track, and that
clearly distinguishes deterministic overlay replay, which does not
re-run the model, from checkpoint internal validation
(\hyperref[fig-protocol]{}). Third, results from eight external
evaluations of four checkpoints together with internal validation
(\hyperref[tab-val-ade]{}--\hyperref[tab-integrity]{}). Fourth, the
specification of the same-scene controlled experiment required for a
causal claim about map and route effects, plus an audit table
enumerating the evidence behind each claim (appendix).

\subsection{Related Work}\label{related-work}

\subsubsection{Unified End-to-End
Driving}\label{unified-end-to-end-driving}

UniAD \hyperref[ref-uniad]{}introduced a planning-oriented design in
which detection, tracking, mapping, motion forecasting, occupancy
prediction, and planning are linked through query interaction so that
every task contributes to the final plan. VAD
\hyperref[ref-vad]{}replaced the scene with a vectorized representation
to reduce computation, and TransFuser
\hyperref[ref-transfuser]{}\hyperref[ref-transfuser-cvpr]{}fused camera
and LiDAR features with a transformer and evaluated the result in
closed-loop CARLA. AutoE2E does not reproduce this complete task chain.
It inherits the idea of organizing a shared BEV representation toward
planning, while narrowing its scope to temporal camera BEV, explicit
navigation context, auxiliary route reconstruction, and direct ego
planning. We do not compare against UniAD on a common benchmark and make
no claim of superiority or inferiority.

\subsubsection{Temporal BEV Encoders}\label{temporal-bev-encoders}

BEVFormer \hyperref[ref-bevformer]{}aggregates multi-camera features
into BEV spatiotemporally using learnable BEV queries and deformable
attention derived from Deformable DETR \hyperref[ref-deformable]{}.
BEVFormer v2 \hyperref[ref-bevformerv2]{}adapts modern image backbones
to BEV perception through perspective supervision, and releases an
official ResNet-50 checkpoint together with a temporal configuration
that fuses multi-frame BEV by convolution. The frozen camera pathway in
this paper starts from the official R50 T8 checkpoint, resizes the BEV
queries and positional embeddings to a 300×200 grid, and loads the
weights with the detection head excluded.

\subsubsection{Map- and Route-Conditioned
Planning}\label{map--and-route-conditioned-planning}

ChauffeurNet \hyperref[ref-chauffeurnet]{}rendered the road map and the
intended route as rasters and used them as inputs for imitation
learning. VectorNet \hyperref[ref-vectornet]{}and LaneGCN
\hyperref[ref-lanegcn]{}encoded HD maps as vectors and lane graphs,
showing that map structure governs motion-forecasting accuracy. The
UniAD-family baselines evaluated on the KITScenes
\hyperref[ref-kitscenes]{}E2E benchmark are conditioned on discrete
navigation commands (turn left, turn right, go straight), whereas the
route used here is a raster consisting of a lane-sequence corridor over
a Lanelet2 \hyperref[ref-lanelet2]{}map together with a destination
marker. Isolating the effect of route conditioning on planning requires
a same-scene controlled experiment, whose specification we give in the
appendix.

\subsubsection{Evaluation Protocols}\label{evaluation-protocols}

Open-loop planning evaluation on nuScenes has repeatedly been shown to
admit shortcuts by which high scores can be obtained from the ego state
alone \hyperref[ref-admlp]{}\hyperref[ref-egostatus]{}. Closed-loop
evaluation in nuPlan \hyperref[ref-nuplan]{}, non-reactive simulation in
NAVSIM \hyperref[ref-navsim]{}, the finding that rule-based planners can
rival learned planners \hyperref[ref-pdm]{}, and conditioning-signal
shortcut biases \hyperref[ref-hidden]{}all indicate that the definition
of the evaluation condition itself shapes the conclusions. This paper
remains in the open-loop setting, but adopts three design principles:
align the evaluation inputs with the deployment contract, distinguish
deterministic replay that does not re-run the model from internal
validation, and do not read differences between distinct scene
populations as causal effects.

\subsection{Problem Setting and
Requirements}\label{problem-setting-and-requirements}

\subsubsection{Deployment Contract}\label{deployment-contract}

At every frame, the deployed AutoE2E policy receives the following:
six-view camera images (the current frame plus 7 history frames at 0.5 s
intervals, each 512×512), a 1024×1024 front camera image for the current
frame, ego-centered HD map rasters (14 channels) and a selected-route
raster (2 channels), and 6.4 s of ego-motion history at 10 Hz. The
output is 6.4 s of acceleration and curvature at 10 Hz. The primary
evaluation should therefore be conducted under the condition in which
all of these inputs are available. At the same time, because situations
in which the map or the route is missing can occur in operation, the
behavior obtained from cameras alone is worth measuring as a separate
track. The two tracks answer different questions.

\subsubsection{Principle of Evaluation
Separation}\label{principle-of-evaluation-separation}

(1) The primary benchmark is KITScenes Val v3.5, where Camera + HD Map +
Route are all available, and the numbers are computed over the 117
scenes and 11,035 samples that satisfy strict identity with the official
Map and Route evaluation. (2) Because the KITScenes Test v1.0
configuration does not release maps, it is evaluated as a camera-only
robustness and transfer track over all 206 scenes and 23,690 samples.
(3) Since Val and Test have different scene populations, the
Test-versus-Val difference is not interpreted as a causal effect of the
map and route. (4) The external evaluation is a replay of the published
deterministic control overlays; it does not re-run the model and
contains no route counterfactual inputs, input gradients, or
route-reconstruction outputs. (5) Checkpoint internal validation is a
separate protocol and is not mixed into the same ranking table.

{\def\LTcaptype{none} % do not increment counter
\begingroup\scriptsize
\begin{longtable}[]{@{}llllllll@{}}
\toprule\noalign{}
System & Camera-only primary evaluation & HD map input & Route input &
Temporal BEV & Unified detection / tracking / occupancy & Closed-loop
evaluation & Evaluation data \\
\midrule\noalign{}
\endhead
\bottomrule\noalign{}
\endlastfoot
UniAD\hyperref[ref-uniad]{} & $\circ$ & --- (online map estimation) &
Navigation command & $\circ$ & $\circ$ & --- & nuScenes \\
VAD\hyperref[ref-vad]{} & $\circ$ & --- (vector map estimation) & Navigation
command & $\circ$ & Partial & $\circ$ (CARLA) & nuScenes / CARLA \\
TransFuser\hyperref[ref-transfuser]{} & --- (camera + LiDAR) & --- &
Target point & --- & Auxiliary tasks & $\circ$ (CARLA) & CARLA \\
AutoE2E (this paper) & --- (secondary track) & $\circ$ (14-ch raster) & $\circ$
(2-ch raster) & $\circ$ (T8, 3.5 s) & --- (BEV segmentation head disabled) &
--- & KITScenes Val / Test, nuPlan internal \\
\end{longtable}
\endgroup
}

\subsection{Method}\label{method}

The description in this section is based on facts confirmed by reading
the source code of the evaluated branch and by instantiating the model
under the same configuration. Symbols in the equations correspond to
\hyperref[fig-arch]{}.

\subsubsection{Temporal Camera BEV}\label{temporal-camera-bev}

The camera pathway builds on the official BEVFormer V2 R50 T8 checkpoint
\hyperref[ref-bevformerv2]{}. A four-level, 256-channel feature pyramid
is formed from the last three stages of ResNet-50
\hyperref[ref-resnet]{}, and a six-layer BEVFormer encoder updates
300×200 learned BEV queries (256-dimensional, 15.36 M parameters)
(\hyperref[fig-encoder]{}). The BEV grid covers X from −60 m to 120 m
and Y from −60 m to 60 m in ego coordinates, with 0.6 m cells
(\hyperref[fig-geometry]{}). Each layer consists of deformable
self-attention over two queues obtained by duplicating the current BEV
(8 heads, 4 points each), multi-scale spatial cross-attention that maps
reference points at four pillar heights (z ∈ {[}−5, 3{]} m) into the
images through a calibrated projection operator (8 heads, 4 levels, 8
points each), a 512-dimensional FFN, and three LayerNorms. The
projection operator implements pinhole and fisheye (F-theta) models, so
this evaluation is not calibration-free.

\begin{figure}
\centering
\pandocbounded{\includegraphics[keepaspectratio,width=\linewidth]{figures_pdf/fig_geometry.pdf}}
\caption{BEV spatial extents. The camera BEV latent (300×200, 0.6 m) and
the map and route rasters fed to the model (450×300, 0.4 m) cover the
same extent (X: −60 to 120 m, Y: −60 to 60 m). The dotted outline is the
published 256×256, 1.0 m geometry audited in KITScenes v3, not the
raster the model receives in this evaluation.}\label{fig-geometry}
\end{figure}

In the T8 temporal configuration, seven history frames at 0.5 s
intervals (t−3.5 s to t−0.5 s) and the current frame are passed through
the same backbone and encoder, with gradients detached for the history
BEVs. The eight BEVs are concatenated into 2048 channels, passed through
a three-block ResNet-style fusion module with 512 intermediate channels,
and projected to 256 channels by a linear layer to obtain B\_img (Eq.
\hyperref[eq-temporal]{}). The camera pathway therefore observes 3.5 s
at 2 Hz, a temporal span different from the 6.4 s ego-motion history
described below (\hyperref[fig-temporal]{}).

\protect\phantomsection\label{eq-temporal}
\(B_{img} = T\left( B_{- 7},B_{- 6},\ldots,B_{- 1},B_{0} \right)\)

\begin{figure}
\centering
\pandocbounded{\includegraphics[keepaspectratio,width=\linewidth]{figures_pdf/fig02_temporal_contract_v2.pdf}}
\caption{Temporal contract. The camera branch uses 8 frames at 0.5 s
intervals (3.5 s), the ego-motion history uses 64 steps at 10 Hz (6.4
s), and the output is 64 steps at 10 Hz (6.4 s). Under the KITScenes
benchmark protocol, only 40 history steps (4.0 s, the remainder
zero-padded) and 50 future steps (5.0 s) are real data, so the effective
horizon coverage of external replay is 50/64 = 78.125\%, and 6.4 s
ADE/FDE cannot be obtained from external replay.}\label{fig-temporal}
\end{figure}

\begin{figure}
\centering
\pandocbounded{\includegraphics[keepaspectratio,width=\linewidth]{figures_pdf/fig03_encoder_layer_v2.pdf}}
\caption{BEVFormer V2 encoder layer (left, six layers, frozen) and the
front-camera residual branch (right). The front branch is a frozen copy
of the last layer\textquotesingle s cross-attention that computes only a
content-dependent delta from the 1024×1024 current front image,
multiplies it by the tanh of a zero-initialized 256-dimensional gate,
and adds it to the current BEV. It is exactly the identity at
initialization, and this gate is the only camera-side parameter that
remains trainable after freezing. History frames use 512×512
only.}\label{fig-encoder}
\end{figure}

\subsubsection{Asymmetric-resolution front
branch}\label{asymmetric-resolution-front-branch}

The current-frame front camera is additionally processed at 1024×1024.
The usual 512 px front view remains in the all-view encoder, while a
branch consisting of a frozen copy of the last encoder
layer\textquotesingle s cross-attention computes a content-dependent
delta δ from the high-resolution front features and adds it as B\_0 =
B\_enc + tanh(g) ⊙ δ. Since g is 256-dimensional and zero-initialized,
the output of the pre-trained camera pathway is unchanged at the start
of training. These 256 parameters are the only camera-side parameters
excluded from freezing. No 1024 px images are used for history frames.

\subsubsection{Navigation rasters}\label{navigation-rasters}

The HD map and the route are represented as ego-centric semantic
rasters. The rasters fed to the evaluated checkpoints are 450×300 cells
at 0.4 m/px, covering the same extent as the camera BEV (X: −60 to 120
m, Y: −60 to 60 m). The semantics of the 14 map-raster channels and the
2 route-raster channels are given in \hyperref[tab-channels]{}, and
their rendering on a real KITScenes Val v3.5 sample in
\hyperref[fig-channels]{}. Binary channels store 1, direction channels
store (sin θ + 1)/2 and (cos θ + 1)/2, and road level stores
(clip(level, −8, 8) + 8)/16, so all values lie in {[}0, 1{]}. The route
corridor is drawn from the sequence of lane centerlines with a width of
3.5 m (truncated 10 m behind the ego vehicle), and the destination
marker is a circle of radius 2.0 m. Map validity and route validity are
explicit per-sample gates, and invalid rasters are multiplied by zero.
Route conditioning can be disabled either in the checkpoint or at
evaluation time, enabling controlled comparisons on the same scene.

{\def\LTcaptype{none} % do not increment counter
\begingroup\scriptsize
\begin{longtable}[]{@{}llllll@{}}
\toprule\noalign{}
ch & Name & Source primitive (Lanelet2 \hyperref[ref-lanelet2]{}) &
Value & Mean occupancy {[}\%{]} & Non-empty / 30 \\
\midrule\noalign{}
\endhead
\bottomrule\noalign{}
\endlastfoot
0 & Drivable area & Lanelet polygons (excluding crosswalks) & binary &
15.55 & 30 \\
1 & Lane boundary & Left and right boundary lines (width 1.0 m) & binary
& 8.29 & 30 \\
2 & Lane centerline & Centerline (width 1.0 m) & binary & 5.52 & 30 \\
3 & Intersection & Lanelet polygons with a turn\_direction attribute &
binary & 1.46 & 28 \\
4 & Crosswalk & Polygons of subtype crosswalk & binary & 0.96 & 24 \\
5 & Stop line & Stop-line polylines (width 1.0 m) & binary & 0.18 &
27 \\
6 & Static traffic light & Positions of traffic\_light regulatory
elements (radius 1.0 m) & binary & 0.07 & 26 \\
7 & Heading sin & Centerline orientation (width 3.5 m) & (sin+1)/2 &
14.72 & 30 \\
8 & Heading cos & Same as above & (cos+1)/2 & 14.72 & 30 \\
9 & Heading valid & Same as above & binary & 14.72 & 30 \\
10 & Known map area & Map boundary polygon & binary & 96.98 & 30 \\
11 & Road level & layer / level attributes & (level+8)/16 & 0.00 & 0 \\
12 & Road level valid & Same as above & binary & 0.00 & 0 \\
13 & Overlapping-level ambiguity & Overwrites by a different level, or
mismatch with the route level & binary & 0.00 & 0 \\
R0 & Selected route corridor & Centerlines of the route lane sequence
(width 3.5 m, truncated 10 m behind) & binary & 1.45 & 28 \\
R1 & Destination marker & Route endpoint (radius 2.0 m) & binary & 0.05
& 26 \\
\end{longtable}
\endgroup
}

\begin{figure}
\centering
\pandocbounded{\includegraphics[keepaspectratio,width=\linewidth]{figures_pdf/fig_channels_real.pdf}}
\caption{The 14 map channels and 2 route channels for one KITScenes Val
v3.5 sample (450×300, 0.4 m/px). These are real data taken from the
dataset artifacts of the public dashboard; red marks the ego position
and the recorded 5 s trajectory. The sample selection rule is given in
Appendix A.5. Across the 30 scanned samples, the road-level channels
11--13 were all empty.}\label{fig-channels}
\end{figure}

The KITScenes route is the lane sequence obtained by fitting the ego
trajectory over the whole scene to the Lanelet2 map, and the destination
is the ego position at the end of the scene (\hyperref[fig-route]{}).
The future ego points themselves are not drawn, but the route is a
retrospectively determined sequence of already-driven lanes, that is, it
equals the intent the driver actually followed. Because this fact bears
directly on how route conditioning should be interpreted, we return to
it in Section 7.

\begin{figure}
\centering
\pandocbounded{\includegraphics[keepaspectratio,width=\linewidth]{figures_pdf/fig14_route_provenance.pdf}}
\caption{Generation of the map and route rasters in KITScenes.
Primitives are extracted from the Lanelet2 map, and the ego trajectory
over the whole scene is fitted to a lane sequence using HMM-like
transition costs to yield the route. Rasters are drawn at anchor poses
spaced 500 ms apart and warped to each sample pose by an SE(2)
transform. The route is not the output of a live route planner but a
retrospective lane-level driving intent.}\label{fig-route}
\end{figure}

\subsubsection{Separate encoding and gated
addition}\label{separate-encoding-and-gated-addition}

The static map and the selected route have separate encoders
(\hyperref[fig-navfusion]{}). The map encoder E\_map is a fully
convolutional network in which a 5×5 convolutional stem (96 channels)
and a context path reduced to strides 2 and 4 by depthwise convolutions
(192 channels) are aligned to 256 channels by 1×1 projections and added,
with the input bilinearly resampled to 300×200. The route encoder
E\_route is a lightweight convolutional network with 96 hidden channels,
designed so that the thin route corridor and the destination features
are not lost to coarse pooling. The route contribution is controlled by
a 256-dimensional per-channel gate and added to the encoded map (Eq.
\hyperref[eq-nav]{}). The gate g is zero-initialized, so the weight at
the start of training is 0.5.

\protect\phantomsection\label{eq-nav}
\(B_{nav} = E_{map}{(M)} + \sigma{(g)} \odot E_{route}{(R)}\)

\begin{figure}
\centering
\pandocbounded{\includegraphics[keepaspectratio,width=\linewidth]{figures_pdf/fig06_navigation_fusion.pdf}}
\caption{Internal structure of the separate navigation encoding and the
deformable navigation fusion. Left: map encoder E\_map (0.167 M) and
route encoder E\_route (0.078 M), the per-channel gate, and the
addition. Right: deformable cross-attention that uses each cell of
B\_img as a query and samples 8 points on B\_nav (8 heads, 256 channels,
0.350 M). The output projection and the final linear layer of the FFN
are zero-initialized, so the fusion is exactly the identity at the start
of training. The gated route contribution feeds the route reconstruction
head (0.017 M).}\label{fig-navfusion}
\end{figure}

\subsubsection{Deformable navigation
fusion}\label{deformable-navigation-fusion}

The navigation BEV is fused into the camera BEV by deformable
cross-attention \hyperref[ref-deformable]{}(Eq. \hyperref[eq-fuse]{}).
Each image-BEV cell acts as a query, predicts 8 offsets from its own
position as the reference point, bilinearly samples B\_nav, and
aggregates the samples with 8-head softmax weights. This design reduces
what would be O(N²) for dense cross-attention over 60,000 cells to
O(NK). Because the output projection and the final linear layer of the
FFN are zero-initialized, the fusion is residual and does not alter the
representation of the pre-trained camera pathway at the start of
training. This acts as a stabilization mechanism when attaching new
navigation features to a frozen visual BEV representation.

\protect\phantomsection\label{eq-fuse}
\(B_{fused} = B_{img} + F_{deform}\left( B_{img},B_{nav} \right)\)

\subsubsection{Ego-motion history and GRU control
planner}\label{ego-motion-history-and-gru-control-planner}

The ego-motion history h\_ego is the flattened 256-dimensional vector of
{[}speed, acceleration, yaw rate, curvature{]} over 64 steps at 10 Hz,
with speed normalized by 33 m/s and acceleration by 8 m/s². In KITScenes
it is derived from finite differences of the pose sequence, and
curvature is the yaw rate divided by the speed lower-bounded at 0.1 m/s.
The model API retains an 896-dimensional visual history input h\_vis,
but the evaluated checkpoints use enable\_world\_model = false and
temporal\_memory\_mode = no\_memory, so h\_vis is fed with zeros. JEPA
training, a learned world model, and a learned 6.4 s visual history
encoder are absent from these checkpoints.

\begin{figure}
\centering
\pandocbounded{\includegraphics[keepaspectratio,width=\linewidth]{figures_pdf/fig07_gru_planner.pdf}}
\caption{Deterministic GRU control planner (0.835 M) and unicycle
integration. The initial state is the sum of linear projections of the
ego-motion history and the visual history. At each step, sampling
locations are predicted from the sum of the recurrent state (gradients
detached) and a learned ego query, 16 points on the fused BEV are
aggregated to update the GRU, and acceleration and curvature are
emitted. This procedure is repeated for 64 steps.}\label{fig-planner}
\end{figure}

The planner is a deterministic GRU (\hyperref[fig-planner]{}). The
Bezier and flow-matching planners present in the repository are not used
in the evaluated checkpoints. The initial context c\_0 is given by Eq.
\hyperref[eq-ctx]{}; at each future step, 16 points on the fused BEV are
attended to deformably (Eq. \hyperref[eq-attn]{}), and the GRU updates
its state and emits acceleration and curvature (Eqs.
\hyperref[eq-gru]{}, \hyperref[eq-out]{}). This is repeated for 64
steps.

\protect\phantomsection\label{eq-ctx}
\(c_{0} = P_{ego}{(h_{ego})} + P_{vis}{(h_{vis})}\)

\protect\phantomsection\label{eq-attn}
\(a_{t} = A_{deform}\left( q_{t},B_{fused} \right),\quad q_{t} = {sg}{(h_{t - 1})} + e\)

\protect\phantomsection\label{eq-gru}
\(h_{t} = {GRU}\left( a_{t},h_{t - 1} \right)\)

\protect\phantomsection\label{eq-out}
\(u_{t} = W_{u}h_{t} = \left( a_{t}^{acc},\kappa_{t} \right)\)

\subsubsection{Unicycle integration}\label{unicycle-integration}

Evaluation and the training loss convert the control sequence into
ego-frame XY with a semi-implicit unicycle model. The time step is Δt =
0.1 s, the initial speed is the current speed v\_0, and speed is clipped
to be non-negative (Eq. \hyperref[eq-roll]{}). The integration is
carried out in float32. No implementation-specific clipping values exist
other than the one on speed.

\protect\phantomsection\label{eq-roll}
\(v_{t} = {\max}\left( v_{t - 1} + a_{t}^{acc}{\Delta}t,0 \right),\;\theta_{t} = \sum_{i \leq t}v_{i}\kappa_{i}{\Delta}t,\; x_{t} = \sum_{i \leq t}v_{i}{\cos\,}\theta_{i}{\Delta}t,\; y_{t} = \sum_{i \leq t}v_{i}{\sin\,}\theta_{i}{\Delta}t\)

\subsubsection{Auxiliary tasks and training
objectives}\label{auxiliary-tasks-and-training-objectives}

The route reconstruction head decodes two channels of logits at 450×300
from the gated route contribution, supervised by the mean of BCE and
soft Dice for the corridor and by a heatmap focal loss (weight 0.25) for
the destination. The trajectory loss is the Smooth-L1 ($\beta$ = 1 m) between
the integrated XY and the recorded trajectory, averaged over valid
steps. All four reported checkpoints were trained with a trajectory
weight of 1.0, a route reconstruction weight of 1.0, and a BEV
segmentation weight of 0.0. A BEV segmentation head (8 classes) exists
in the checkpoints, but the reported checkpoints are not jointly
optimized with BEV segmentation. The perception branch and the
world-model branch are also disabled.

\subsubsection{Parameter budget}\label{parameter-budget}

\hyperref[tab-params]{}and \hyperref[fig-params]{}report parameter
counts obtained by instantiating the model in the same configuration (6
views). Of the 79,906,522 total parameters, 78,293,196 (97.98\%) belong
to the camera-BEV pathway. In total-capacity terms, the model is
therefore almost entirely a pretrained, frozen visual representation: T8
temporal fusion accounts for 38.56\%, ResNet-50 for 29.42\%, the learned
BEV queries for 19.22\%, the BEVFormer encoder layers for 6.18\%, and
the FPN for 4.10\%. The 1,448,582 parameters updated during trajectory
training (1.81\% of the model) are allocated to the GRU planner
(835,380; 57.67\% of trainable parameters), deformable navigation fusion
(350,288; 24.18\%), Map/Route encoders (245,440; 16.94\%),
route-reconstruction head (17,218; 1.19\%), and front residual gate
(256; 0.02\%). Thus, the system does not train all 79.91 M parameters
for planning; it trains a 1.45 M task-specific layer on top of a large
frozen camera-BEV representation. The difference of 512 from the
79,907,034 parameters recorded for the repository\textquotesingle s
8-view benchmark configuration matches two additional camera embeddings.

{\def\LTcaptype{none} % do not increment counter
\begingroup\scriptsize
\begin{longtable}[]{@{}lll@{}}
\toprule\noalign{}
Component & Parameters & State \\
\midrule\noalign{}
\endhead
\bottomrule\noalign{}
\endlastfoot
ResNet-50 backbone & 23,508,032 & frozen \\
Feature pyramid (4 levels) & 3,278,592 & frozen \\
BEV queries 300×200×256 & 15,360,000 & frozen \\
Row, column, level, and camera embeddings & 66,560 & frozen \\
Pseudo-projection matrix (unused in calibrated evaluation) & 12 &
frozen \\
Encoder, 6 layers (823,488 per layer) & 4,940,928 & frozen \\
Front cross-attention (copy) & 328,960 & frozen \\
Front residual gate & 256 & trained \\
T8 temporal fusion & 30,809,856 & frozen \\
Map encoder E\_map & 167,200 & trained \\
Route encoder E\_route & 77,984 & trained \\
Route gate & 256 & trained \\
Deformable navigation fusion & 350,288 & trained \\
GRU control planner & 835,380 & trained \\
Route reconstruction head & 17,218 & trained \\
BEV segmentation head (weight 0.0) & 165,000 & frozen \\
Total / trainable & 79,906,522 / 1,448,582 & --- \\
\end{longtable}
\endgroup
}

\begin{figure}
\centering
\pandocbounded{\includegraphics[keepaspectratio,width=\linewidth]{figures_pdf/fig_param_decomposition_v2.pdf}}
\caption{Parameter allocation of the evaluated six-view configuration.
Panel (a) partitions all 79,906,522 parameters; the frozen camera-BEV
pathway contains 78.29 M (97.98\%). Panel (b) partitions the 1,448,582
trainable parameters among task-specific planning and navigation
modules. The GRU planner contains 57.67\% of trainable parameters,
deformable navigation fusion 24.18\%, and the Map/Route encoders plus
route gate 16.94\%.}\label{fig-params}
\end{figure}

\subsection{Datasets and Evaluation
Protocol}\label{datasets-and-evaluation-protocol}

\subsubsection{Training Data and
Provenance}\label{training-data-and-provenance}

All four trajectory checkpoints were trained with 8 distributed GPU
workers, training seed 149, a frozen BEVFormer camera branch, trajectory
weight 1.0, route-reconstruction weight 1.0, and BEV segmentation weight
0.0 (\hyperref[tab-config]{}, \hyperref[fig-protocol]{}). The nuPlan
\hyperref[ref-nuplan]{}stage checkpoints (Epoch 4 and Epoch 5) use a
learning rate of 1e-4 and an internal validation set of 1,024 samples.
KITScenes \hyperref[ref-kitscenes]{}fine-tuning uses the frozen nuPlan
Epoch 5 as its parent with a learning rate of 3e-5. The split manifest
identifies the source as the official KITScenes training split: 129
empty scenes are excluded from 533 available train scenes, leaving 404
eligible scenes and 42,667 samples. A scene-level 90\%/10\% split yields
38,847 samples from 364 scenes for training and 3,820 samples from 40
scenes for frozen internal validation (BF16 validation precision).
Because these 404 scenes belong to the official training split, their
scene IDs are disjoint from the official 117-scene Val and 206-scene
Test populations. KITScenes Epoch 5 was the best checkpoint retained at
that point; Epoch 7 was not marked best. The selection score used by
that registry flag is not documented in the evidence available to this
paper. The registration metadata records eval\_gate\_pass = false for
all four checkpoints, and this paper does not claim that a quality gate
passed (Appendix A.3).

\begin{minipage}[t]{0.45\linewidth}
{\def\LTcaptype{none} % do not increment counter
\begingroup\scriptsize
\begin{longtable}[]{@{}lllllll@{}}
\toprule\noalign{}
Checkpoint & SHA-256 prefix & Training stage & Learning rate & Parent &
Internal validation samples & Notes \\
\midrule\noalign{}
\endhead
\bottomrule\noalign{}
\endlastfoot
nuPlan Epoch 4 & ed00e072471a & nuPlan trajectory and route (frozen
camera BEV) & 1e-4 & Official BEVFormer V2 R50 T8 & 1,024 & Registered
as best on nuPlan \\
nuPlan Epoch 5 & ca8b43d7a777 & Same as above & 1e-4 & Same as above &
1,024 & Parent of KITScenes fine-tuning \\
KITScenes Epoch 5 & 120a21639d97 & KITScenes fine-tuning & 3e-5 & nuPlan
Epoch 5 (frozen) & 3,820 & Retained best \\
KITScenes Epoch 7 & a1e6b1621018 & KITScenes fine-tuning & 3e-5 & nuPlan
Epoch 5 (frozen) & 3,820 & No best designation \\
\end{longtable}
\endgroup
}
\end{minipage}
\hfill
\begin{minipage}[t]{0.45\linewidth}
Common to all: 8 GPU workers, seed 149, trajectory weight 1.0,
route-reconstruction weight 1.0, BEV segmentation weight 0.0,
eval\_gate\_pass = false.
\end{minipage}

\begin{figure}
\centering
\pandocbounded{\includegraphics[keepaspectratio,width=\linewidth]{figures_pdf/fig08_lineage_protocol.pdf}}
\caption{Top: checkpoint lineage. Starting from the official BEVFormer
V2 R50 T8, the nuPlan trajectory stage (Epoch 4 / 5), then KITScenes
fine-tuning with the frozen Epoch 5 as parent (Epoch 5 / 7). Bottom: the
two evaluation protocols. Protocol A is checkpoint internal validation
that runs the model (6.4 s); Protocol B is replay of the released
deterministic control overlays (up to 5 s). The two are never mixed in a
single ranking table. Val and Test are different scene
populations.}\label{fig-protocol}
\end{figure}

\subsubsection{Evaluation datasets}\label{evaluation-datasets}

The primary evaluation is KITScenes Val v3.5 under a Camera + HD Map +
oracle post-hoc Route input condition. The public dashboard displays 140
shard artifacts and 13,525 samples, whereas numerical evaluation uses
the 117 official Val scenes and 11,035 samples satisfying the stricter
Map/Route evaluation identity. The available report does not contain a
per-reason breakdown for the 2,490 excluded samples; selection bias from
this identity filter therefore cannot be ruled out. The map-valid and
route-valid fractions of the 11,035-sample population were not retained
in the replay report. The secondary evaluation is KITScenes Test v1.0
over all 206 scenes and 23,690 samples. Its manifests mark Map and Route
unavailable; the loader supplies zero-valued rasters with validity gates
false before the encoders. Because trained GroupNorm affine terms may
respond to a zero tensor, this condition is more precisely an
invalid-navigation-input condition than a structural removal of the
navigation modules. The official KITScenes E2E protocol uses 200
nine-second windows from validation data, with 4 s of observation and up
to 5 s of future trajectory \hyperref[ref-kitscenes]{}. Our replay
approximates its 40-history/50-future temporal contract but applies it
to all samples admitted by the official Val/Test identity. Val and Test
are different scene populations and are never treated as a paired
ablation.

\subsubsection{Protocol A: checkpoint internal
validation}\label{protocol-a-checkpoint-internal-validation}

The model is run with the same code as the training job to compute
ADE/FDE at 6.4 s (64 steps), route-reconstruction IoU, and open-loop
route, drivable-area, and comfort metrics. Route corridor compliance
rate is the fraction of steps at which all four corners of the 4.8 m ×
2.0 m ego footprint lie inside a 3.5 m wide route corridor, and success
rate is the fraction of samples that comply at all 64 steps.
Drivable-area compliance and success rates adopt the same definitions
with respect to the drivable-area mask. The projected-endpoint
arc-length proxy is the arc length obtained by projecting the predicted
endpoint onto the logged trajectory polyline, and the ratio normalizes
it by the logged trajectory length. For comfort, a sample is counted as
a violation if it exceeds any of the nuPlan thresholds (longitudinal
acceleration −4.05 to 2.40 m/s², lateral acceleration 4.89 m/s², yaw
rate 0.95 rad/s, yaw acceleration 1.93 rad/s², longitudinal jerk 4.13
m/s³, jerk magnitude 8.37 m/s³), and the comfort rate is 1 − violation
rate. Route metrics are computed on the 2,911 samples with a valid
route.

\subsubsection{Protocol B: external deterministic overlay
replay}\label{protocol-b-external-deterministic-overlay-replay}

The external evaluation replays the published control overlays for each
checkpoint (64 acceleration/curvature pairs and an initial speed per
sample) through Eq. \hyperref[eq-roll]{}and compares the result with the
logged trajectory; it does not re-run the model. Registry roles state
that the Val overlays were generated by inference with Camera, HD Map
and Route, and Test overlays under camera\_only\_missing\_map\_route.
Replay reproduces the controls but cannot independently verify those
inference inputs. ADE\_h is the per-sample mean Euclidean error over
valid steps through h seconds, then averaged across samples; FDE\_h is
the mean error at h seconds. Lateral and longitudinal errors average
\textbar Δy\textbar{} and \textbar Δx\textbar{} over all valid 0--5 s
steps and do not identify a horizon or mechanism. A sample is non-finite
if any predicted control is non-finite. The replay artifacts contain no
route counterfactual, input gradient, or route-reconstruction output.
Under the benchmark window, 40 real egomotion steps (4.0 s) are
left-zero-padded to the 64-step ABI and only 50 future steps (5.0 s)
have valid targets. We therefore report 1, 2, 3 and 5 s metrics, not
external 6.4 s metrics. By contrast, Protocol A uses the training-window
construction with 64 real future targets; its 6.4 s KITScenes results
are a different sample construction and are not ranked with Protocol B.
No kinematic reference baseline or paired uncertainty interval is
available in the completed replay report. Bold and underline in the
tables denote numerical ordering only, not statistical significance.

\subsection{Results}\label{results}

\subsubsection{Primary evaluation: KITScenes Val (Camera + HD Map +
Route)}\label{primary-evaluation-kitscenes-val-camera--hd-map--route}

\hyperref[tab-val-ade]{}and \hyperref[fig-val]{}report ADE/FDE for the
four checkpoints on 11,035 samples. KITScenes Epoch 7 has the lowest
reported means at 1 s in ADE (0.1347 m) and FDE (0.2831 m), while
KITScenes Epoch 5 has the lowest reported means at 2, 3 and 5 s for both
ADE and FDE (5 s: ADE 1.9405 m, FDE 5.5645 m). The two nuPlan
checkpoints are inferior to the KITScenes fine-tuned ones at every
horizon, with 5 s FDE of 6.2--6.4 m.

{\def\LTcaptype{none} % do not increment counter
\begingroup\scriptsize
\begin{longtable}[]{@{}lllllllll@{}}
\toprule\noalign{}
Model & ADE 1 s & ADE 2 s & ADE 3 s & ADE 5 s & FDE 1 s & FDE 2 s & FDE
3 s & FDE 5 s \\
\midrule\noalign{}
\endhead
\bottomrule\noalign{}
\endlastfoot
KITScenes Epoch 5 & 0.1472 & 0.3729 & 0.7449 & 1.9405 & 0.2850 & 0.9239
& 2.0196 & 5.5645 \\
KITScenes Epoch 7 & 0.1347 & 0.3886 & 0.8035 & 2.0952 & 0.2831 & 1.0101
& 2.2127 & 5.9509 \\
nuPlan Epoch 5 & 0.1951 & 0.4723 & 0.9147 & 2.2951 & 0.3779 & 1.1331 &
2.4226 & 6.3885 \\
nuPlan Epoch 4 & 0.2030 & 0.5018 & 0.9503 & 2.2919 & 0.4066 & 1.1935 &
2.4508 & 6.2272 \\
\end{longtable}
\endgroup
}

\begin{figure}
\centering
\pandocbounded{\includegraphics[keepaspectratio,width=\linewidth]{figures_pdf/fig_results_val.pdf}}
\caption{ADE (left) and FDE (right) across horizons on KITScenes Val
v3.5 (Camera + HD Map + Route; n = 11,035). Solid lines are the
KITScenes fine-tuned checkpoints, dashed lines the nuPlan-trained ones.
The 5 s values are annotated.}\label{fig-val}
\end{figure}

\subsubsection{Secondary evaluation: KITScenes Test (Camera
only)}\label{secondary-evaluation-kitscenes-test-camera-only}

\hyperref[tab-test-ade]{}and \hyperref[fig-test]{}report results on the
camera-only track with 23,690 samples. KITScenes Epoch 5 has the lowest
reported means at 1 s (ADE 0.1463 m, FDE 0.3107 m) and 5 s (ADE 2.1042
m, FDE 5.9422 m), whereas nuPlan Epoch 5 is lowest at 2 s and 3 s.
KITScenes Epoch 7 has the largest errors at 2, 3 and 5 s; at 1 s, nuPlan
Epoch 4 has the largest ADE and FDE. Because Val and Test have different
scene populations, placing this table beside
\hyperref[tab-val-ade]{}does not measure the effect of map and route.

{\def\LTcaptype{none} % do not increment counter
\begingroup\scriptsize
\begin{longtable}[]{@{}lllllllll@{}}
\toprule\noalign{}
Model & ADE 1 s & ADE 2 s & ADE 3 s & ADE 5 s & FDE 1 s & FDE 2 s & FDE
3 s & FDE 5 s \\
\midrule\noalign{}
\endhead
\bottomrule\noalign{}
\endlastfoot
KITScenes Epoch 5 & 0.1463 & 0.4121 & 0.8263 & 2.1042 & 0.3107 & 1.0453
& 2.2199 & 5.9422 \\
KITScenes Epoch 7 & 0.1680 & 0.4923 & 0.9831 & 2.4699 & 0.3754 & 1.2512
& 2.6242 & 6.8964 \\
nuPlan Epoch 5 & 0.1602 & 0.3979 & 0.7924 & 2.1120 & 0.3134 & 0.9727 &
2.1621 & 6.1751 \\
nuPlan Epoch 4 & 0.1916 & 0.4817 & 0.9181 & 2.2840 & 0.3899 & 1.1514 &
2.3895 & 6.4266 \\
\end{longtable}
\endgroup
}

\begin{figure}
\centering
\pandocbounded{\includegraphics[keepaspectratio,width=\linewidth]{figures_pdf/fig_results_test.pdf}}
\caption{ADE and FDE on KITScenes Test v1.0 (Camera only; n = 23,690).
This is a different scene population from Val, so differences with
respect to \hyperref[fig-val]{}do not imply a causal effect of map and
route.}\label{fig-test}
\end{figure}

\subsubsection{Lateral and longitudinal
error}\label{lateral-and-longitudinal-error}

\hyperref[tab-latlon]{}and \hyperref[fig-latlon]{}report the mean
absolute lateral error (\textbar Δy\textbar) and longitudinal error
(\textbar Δx\textbar) over all valid steps. On Val, KITScenes Epoch 5 is
smallest in both (0.9844 m, 1.4025 m). On Test, nuPlan Epoch 5 is
smallest laterally (1.0584 m) and KITScenes Epoch 5 longitudinally
(1.4051 m). For every model, mean longitudinal error exceeds mean
lateral error over all valid 0--5 s steps. This aggregate does not
identify a horizon or a causal mechanism, and the scene population has
much greater along-track than cross-track displacement.

{\def\LTcaptype{none} % do not increment counter
\begingroup\scriptsize
\begin{longtable}[]{@{}lllll@{}}
\toprule\noalign{}
\multirow{2}{*}{Model} & \multicolumn{2}{l}{%
KITScenes Val v3.5 (Camera + Map + Route)} & \multicolumn{2}{l@{}}{%
KITScenes Test v1.0 (Camera only)} \\
& Lateral \textbar Δy\textbar{} & Longitudinal \textbar Δx\textbar{} &
Lateral \textbar Δy\textbar{} & Longitudinal \textbar Δx\textbar{} \\
\midrule\noalign{}
\endhead
\bottomrule\noalign{}
\endlastfoot
KITScenes Epoch 5 & 0.9844 & 1.4025 & 1.2044 & 1.4051 \\
KITScenes Epoch 7 & 1.0220 & 1.5659 & 1.4730 & 1.6598 \\
nuPlan Epoch 5 & 1.1911 & 1.6633 & 1.0584 & 1.5496 \\
nuPlan Epoch 4 & 1.2082 & 1.6448 & 1.2167 & 1.6287 \\
\end{longtable}
\endgroup
}

\begin{figure}
\centering
\pandocbounded{\includegraphics[keepaspectratio,width=\linewidth]{figures_pdf/fig_results_latlon.pdf}}
\caption{Mean absolute lateral error (light) and longitudinal error
(hatched). Left: Val (Camera + Map + Route); right: Test (Camera only).
The two panels are evaluations on different
populations.}\label{fig-latlon}
\end{figure}

\subsubsection{Checkpoint internal
validation}\label{checkpoint-internal-validation}

\hyperref[tab-internal]{}reports the internal validation protocol and is
treated as a ranking table separate from the external replay. On nuPlan
internal validation (1,024 samples) the nuPlan checkpoints reach ADE of
about 1.26 m and FDE of about 3.85--3.90 m at 6.4 s, whereas on the
KITScenes frozen validation (3,820 samples) the KITScenes checkpoints
reach ADE 2.63--2.75 m and FDE 7.56--7.86 m. Since these are numbers on
different datasets, the nuPlan and KITScenes rows are likewise not
compared with each other. Route-reconstruction IoU is above 0.99 in all
cases, but because the reconstruction head receives route-derived
features, this metric mainly verifies that route information is
preserved in the representation and does not mean that the vehicle drove
along the route.

{\def\LTcaptype{none} % do not increment counter
\begingroup\scriptsize
\begin{longtable}[]{@{}lllll@{}}
\toprule\noalign{}
Model & n & ADE 6.4 s & FDE 6.4 s & Route reconstruction IoU \\
\midrule\noalign{}
\endhead
\bottomrule\noalign{}
\endlastfoot
nuPlan Epoch 4 & 1,024 & 1.2578 & 3.8450 & 0.9912 \\
nuPlan Epoch 5 & 1,024 & 1.2668 & 3.8982 & 0.9980 \\
KITScenes Epoch 5 & 3,820 & 2.6326 & 7.5567 & 0.9991 \\
KITScenes Epoch 7 & 3,820 & 2.7548 & 7.8574 & 0.9990 \\
\end{longtable}
\endgroup
}

For the two KITScenes checkpoints, \hyperref[tab-internal-kit]{}and
\hyperref[fig-tradeoff]{}report the open-loop route, drivable-area, and
comfort metrics. Epoch 7 increases route corridor compliance rate
(0.3959 → 0.4236), route corridor success rate (0.0608 → 0.0763), the
projected-endpoint arc-length proxy (44.7857 → 45.5856 m; ratio 0.8725 →
0.8982), and increases drivable-area compliance rate (0.7758 → 0.7988),
while degrading drivable-area success rate (0.4442 → 0.4139) and comfort
rate (0.7542 → 0.7466), and also degrading ADE/FDE at 6.4 s. In the
breakdown of comfort violations (Appendix A.4), jerk magnitude and yaw
acceleration violations account for the majority, while longitudinal and
lateral acceleration and yaw rate violations are below 2\%.

{\def\LTcaptype{none} % do not increment counter
\begingroup\scriptsize
\begin{longtable}[]{@{}lll@{}}
\toprule\noalign{}
Metric & Epoch 5 & Epoch 7 \\
\midrule\noalign{}
\endhead
\bottomrule\noalign{}
\endlastfoot
Route corridor compliance rate ↑ & 0.3959 & 0.4236 \\
Route corridor success rate ↑ & 0.0608 & 0.0763 \\
Route progress proxy {[}m{]} ↑ & 44.7857 & 45.5856 \\
Route progress proxy ratio ↑ & 0.8725 & 0.8982 \\
Drivable-area compliance rate ↑ & 0.7758 & 0.7988 \\
Drivable-area success rate ↑ & 0.4442 & 0.4139 \\
Comfort rate ↑ & 0.7542 & 0.7466 \\
Comfort violation rate ↓ & 0.2458 & 0.2534 \\
ADE 6.4 s {[}m{]} ↓ & 2.6326 & 2.7548 \\
FDE 6.4 s {[}m{]} ↓ & 7.5567 & 7.8574 \\
\end{longtable}
\endgroup
}

\begin{figure}
\centering
\pandocbounded{\includegraphics[keepaspectratio,width=\linewidth]{figures_pdf/fig_tradeoff_ep5_ep7.pdf}}
\caption{Relative change from KITScenes Epoch 5 to Epoch 7 (internal
validation). Green denotes change in the desirable direction for each
metric, red change in the undesirable direction. The route proxies
improve, while long-horizon trajectory error, drivable-area success
rate, and comfort rate degrade.}\label{fig-tradeoff}
\end{figure}

\subsubsection{Numerical integrity and
coverage}\label{numerical-integrity-and-coverage}

\hyperref[tab-integrity]{}reports the coverage of the external replay.
On Val, 551,750 of 11,035 samples × 64 steps = 706,240 steps are valid;
on Test, 1,184,500 of 23,690 × 64 = 1,516,160, giving a coverage of
78.125\% (50/64) in both cases. Non-finite predictions were 0 in all
eight external evaluations.

{\def\LTcaptype{none} % do not increment counter
\begingroup\scriptsize
\begin{longtable}[]{@{}llllll@{}}
\toprule\noalign{}
Dataset & Samples / model & Valid steps & Total steps & Coverage &
Non-finite predictions \\
\midrule\noalign{}
\endhead
\bottomrule\noalign{}
\endlastfoot
KITScenes Val v3.5 & 11,035 & 551,750 & 706,240 & 78.125\% & 0 \\
KITScenes Test v1.0 & 23,690 & 1,184,500 & 1,516,160 & 78.125\% & 0 \\
\end{longtable}
\endgroup
}

\subsubsection{Qualitative results}\label{qualitative-results}

\hyperref[fig-qual]{}shows six examples in which the released control
overlays are replayed on each sample\textquotesingle s own map and route
raster. Selection followed a fixed rule: sort the shard artifacts by
name, scan the sample at the median index of the first 30 shards, and
take (a) the example with the smallest and (b) the largest 5 s FDE for
Epoch 5, (c) the first left turn, (d) the first right turn, (e) the
first straight approach to an intersection, and (f) the first straight
segment without an intersection (Appendix A.5). No post hoc screening by
error was performed outside this rule. In example (a), Epoch 5 follows
the right turn with an FDE of 0.47 m while Epoch 7 deviates by 5.29 m;
in the roundabout of example (d), Epoch 7 (3.90 m) outperforms Epoch 5
(6.50 m). In example (b), all four models travel farther than the logged
trajectory (FDE 9.9--18.5 m), failing to predict the deceleration. These
are single-sample observations, not statistical conclusions.

\begin{figure}
\centering
\pandocbounded{\includegraphics[keepaspectratio,width=\linewidth]{figures_pdf/fig_qualitative.pdf}}
\caption{Qualitative examples from KITScenes Val v3.5. The top row shows
the model-input front tile of each sample (center crop); the bottom row
shows the logged trajectory (black, 5 s) and the replayed trajectories
of the four checkpoints (first 50 steps) on the map and route raster.
The 5 s FDE is annotated in each panel. The selection rule follows the
main text and Appendix A.5 and involves no post hoc screening by
error.}\label{fig-qual}
\end{figure}

\subsection{Discussion}\label{discussion}

\subsubsection{Epoch 5 versus Epoch 7}\label{epoch-5-versus-epoch-7}

In Protocol B, continuing training from Epoch 5 to Epoch 7 reduced 1 s
error (ADE 0.1472 → 0.1347 m) and degraded accuracy at 2 s and beyond (5
s ADE 1.9405 → 2.0952 m, FDE 5.5645 → 5.9509 m). In Protocol A, the same
additional training increased route-corridor compliance and the censored
projected-endpoint arc-length proxy while increasing 6.4 s error and
lowering comfort. These are direct observations from different protocols
and sample populations, not a joint ranking. For this checkpoint pair,
selecting by one trajectory or route metric would regress another
reported metric. Checkpoint selection should therefore be
multi-objective: The registry marks Epoch 5 as best by an undocumented
selection score; Epoch 7 nevertheless has a lower 1 s mean in Protocol B
and a larger censored projected-endpoint arc-length proxy in Protocol A.
No statistical distinction is claimed.

\subsubsection{Cross-dataset transfer}\label{cross-dataset-transfer}

On the same KITScenes Val scenes, KITScenes fine-tuning outperformed the
nuPlan-trained checkpoints at every horizon (RQ2). The gap between
nuPlan Epoch 4 and Epoch 5 is small, and at 5 s FDE Epoch 4 is slightly
better. On the camera-only Test split the ranking changes: nuPlan Epoch
5 has the lowest reported mean at 2 s and 3 s, whereas KITScenes Epoch 7
is largest within the Test table at 2, 3 and 5 s. KITScenes fine-tuning
was trained under a condition in which the map and the route are
provided, so it is conceivable that it becomes relatively brittle when
they are absent; however, because Test and Val contain different scenes,
geographies, motion distributions, and goal identities, this gap cannot
be quantified as the effect of missing map and route inputs (RQ3).

\subsubsection{Route provenance and route
reconstruction}\label{route-provenance-and-route-reconstruction}

The KITScenes route is a lane sequence obtained by fitting the driven
trajectory of the entire scene to Lanelet2, with the destination set to
the end of the scene (\hyperref[fig-route]{}). This route is equivalent
to the lane-level intent the driver actually chose, and it can be more
informative than a route supplied by a navigation system at deployment
time. The numbers in the primary evaluation should therefore be read as
performance under an oracle route reconstructed from the logged future
trajectory, and they do not cover robustness to route errors or route
replanning. A route-reconstruction IoU of 0.999 verifies that the
reconstruction head can recover the route raster from the post-gate
route contribution, i.e. that the gate does not collapse the route
information; it is not evidence that the policy drove along the route.
Behavioral route adherence is measured by the route-corridor compliance
and success rates, and those values (compliance 0.40--0.42, success
0.06--0.08) are low under the strict definition that a 2.0 m wide
footprint must remain inside a 3.5 m wide corridor at every step.

\subsubsection{What can be claimed about the effect of the map and the
route}\label{what-can-be-claimed-about-the-effect-of-the-map-and-the-route}

The evidence in this paper does not support the claim that the HD map
and the route causally improved planning accuracy (RQ5). There are five
reasons. Val and Test are different scenes. Test, by dataset
construction, has no map and no route. The current report contains no
camera-only baseline on the same scenes as Val. The external overlay
replay does not retain the counterfactual route input. Per-scene paired
errors and bootstrap confidence intervals have not yet been computed.

We instead offer the following design hypotheses. The HD map should
reduce geometric ambiguity in lane topology, intersections, stop lines,
and drivable boundaries. The route should uniquely determine which of
several topologically valid branches the ego vehicle takes. Encoding the
map and the route separately should prevent the sparse chosen route from
being buried in the dense static map features. Residual, gated fusion
should add navigation information without destroying the pretrained
visual BEV representation. These remain hypotheses until controlled
experiments are available.

\subsubsection{Relation to UniAD}\label{relation-to-uniad}

UniAD and AutoE2E both adopt a BEV-centric representation and
planning-oriented feature sharing. UniAD unifies perception, tracking,
mapping, motion forecasting, occupancy, and planning through task
interaction, whereas AutoE2E narrows its scope to temporal camera BEV,
explicit navigation context, auxiliary route reconstruction, and direct
ego planning, and treats externally supplied HD maps and routes as
first-class deployment inputs. The claim of this paper is that
evaluation inputs should match the contract of the deployed policy; we
make no claim of superiority or inferiority relative to UniAD, since no
common benchmark exists. We also do not compare against the zero-shot
UniAD numbers reported in the KITScenes paper (200 samples, 3 s),
because the sample set, horizon, and input conditions differ.

\subsection{Limitations}\label{limitations}

The evaluation in this paper is open-loop only and includes no
closed-loop safety or intervention metrics. No controlled map/route
experiment on identical scenes was conducted. The external replay is
limited to 5 s, so 22\% of the 6.4 s output is not evaluated. Val and
Test are different populations, and the difference between them carries
no causal meaning. The camera encoder is frozen, and the degree to which
a representation trained on nuScenes fits the 6-camera configuration and
0.6 m grid of KITScenes has not been separately evaluated (the BEV
segmentation head is disabled, and occupancy IoU is not computed). The
full UniAD task stack is not included. Confidence intervals are omitted
because per-scene data have not been aggregated. All checkpoint quality
gates are recorded as not passed, and the details of the gate
definitions are not part of the evidence in this paper (Appendix A.3).
The training data are nuPlan and KITScenes from three German cities, so
distribution shifts in geography, season, and map may exist. Because the
route is a post-hoc sequence of already-driven lanes, robustness to
route errors is not measured. Under the KITScenes benchmark protocol,
only 4.0 s of the 6.4 s ego motion history is real data, leaving a
distribution gap with respect to training (6.4 s). The six qualitative
examples and the occupancy rates in \hyperref[tab-channels]{}are
descriptive observations from a scan of 30 samples, not population
statistics.

\subsection{Conclusion}\label{conclusion}

AutoE2E couples separately encoded HD map and oracle post-hoc route
rasters into a frozen temporal camera BEV via a gated residual and
generates controls with a GRU. We evaluated it at scale under matching
camera/map/route modalities, but with 4.0 s of real egomotion history
padded to 6.4 s and only a 5.0 s evaluated output horizon. Under the
Camera + HD Map + Route KITScenes Val protocol, KITScenes Epoch 5 has
the lowest reported Protocol-B means at 2, 3 and 5 s, while Epoch 7 has
the lowest 1 s mean. Separately, Protocol A shows a larger censored
projected-endpoint arc-length proxy for Epoch 7 together with larger 6.4
s error and lower aggregate comfort. The camera-only Test split is
useful for robustness analysis, but because it contains different scenes
it cannot quantify the causal contribution of the navigation inputs.
That claim requires paired Camera / Camera + Map / Camera + Map + Route
evaluation on identical scenes, and this paper specifies that protocol
in the appendix. Deployment-aligned evaluation should retain the
navigation inputs the policy requires, and camera-only evaluation should
be positioned as a complementary robustness track.

\subsection{Acknowledgments}\label{acknowledgments}

We thank the contributors and the community of the Autoware Foundation
Robotaxi Working Group.

\subsection{References}\label{references}

\begin{thebibliography}{99}
\bibitem{ref-uniad}
\begin{otherlanguage}{english}

Y. Hu, J. Yang, L. Chen, K. Li, C. Sima, X. Zhu, S. Chai, S. Du, T. Lin,
W. Wang, L. Lu, X. Jia, Q. Liu, J. Dai, Y. Qiao, and H. Li.
Planning-Oriented Autonomous Driving. In \emph{Proc. IEEE/CVF CVPR}, pp.
17853--17862, 2023. arXiv:2212.10156.

\end{otherlanguage}

\bibitem{ref-bevformer}
\begin{otherlanguage}{english}

Z. Li, W. Wang, H. Li, E. Xie, C. Sima, T. Lu, Y. Qiao, and J. Dai.
BEVFormer: Learning Bird\textquotesingle s-Eye-View Representation from
Multi-camera Images via Spatiotemporal Transformers. In \emph{Computer
Vision -- ECCV 2022}, LNCS 13669, pp. 1--18, Springer, 2022.
doi:10.1007/978-3-031-20077-9\_1.

\end{otherlanguage}

\bibitem{ref-bevformerv2}
\begin{otherlanguage}{english}

C. Yang, Y. Chen, H. Tian, C. Tao, X. Zhu, Z. Zhang, G. Huang, H. Li, Y.
Qiao, L. Lu, J. Zhou, and J. Dai. BEVFormer v2: Adapting Modern Image
Backbones to Bird\textquotesingle s-Eye-View Recognition via Perspective
Supervision. In \emph{Proc. IEEE/CVF CVPR}, pp. 17830--17839, 2023.
arXiv:2211.10439.

\end{otherlanguage}

\bibitem{ref-deformable}
\begin{otherlanguage}{english}

X. Zhu, W. Su, L. Lu, B. Li, X. Wang, and J. Dai. Deformable DETR:
Deformable Transformers for End-to-End Object Detection. In \emph{Proc.
ICLR}, 2021. arXiv:2010.04159.

\end{otherlanguage}

\bibitem{ref-vad}
\begin{otherlanguage}{english}

B. Jiang, S. Chen, Q. Xu, B. Liao, J. Chen, H. Zhou, Q. Zhang, W. Liu,
C. Huang, and X. Wang. VAD: Vectorized Scene Representation for
Efficient Autonomous Driving. In \emph{Proc. IEEE/CVF ICCV}, pp.
8340--8350, 2023. arXiv:2303.12077.

\end{otherlanguage}

\bibitem{ref-transfuser}
\begin{otherlanguage}{english}

K. Chitta, A. Prakash, B. Jaeger, Z. Yu, K. Renz, and A. Geiger.
TransFuser: Imitation with Transformer-Based Sensor Fusion for
Autonomous Driving. \emph{IEEE Trans. Pattern Analysis and Machine
Intelligence}, 45(11):12878--12895, 2023.
doi:10.1109/TPAMI.2022.3200245.

\end{otherlanguage}

\bibitem{ref-transfuser-cvpr}
\begin{otherlanguage}{english}

A. Prakash, K. Chitta, and A. Geiger. Multi-Modal Fusion Transformer for
End-to-End Autonomous Driving. In \emph{Proc. IEEE/CVF CVPR}, pp.
7077--7087, 2021. arXiv:2104.09224.

\end{otherlanguage}

\bibitem{ref-nuplan}
\begin{otherlanguage}{english}

H. Caesar, J. Kabzan, K. S. Tan, W. K. Fong, E. Wolff, A. Lang, L.
Fletcher, O. Beijbom, and S. Omari. nuPlan: A closed-loop ML-based
planning benchmark for autonomous vehicles. \emph{CVPR 2021 Workshop on
Autonomous Driving: Perception, Prediction and Planning (ADP3)}, 2021.
arXiv:2106.11810.

\end{otherlanguage}

\bibitem{ref-nuscenes}
\begin{otherlanguage}{english}

H. Caesar, V. Bankiti, A. H. Lang, S. Vora, V. E. Liong, Q. Xu, A.
Krishnan, Y. Pan, G. Baldan, and O. Beijbom. nuScenes: A Multimodal
Dataset for Autonomous Driving. In \emph{Proc. IEEE/CVF CVPR}, pp.
11621--11631, 2020.

\end{otherlanguage}

\bibitem{ref-kitscenes}
\begin{otherlanguage}{english}

R. Schwarzkopf, F. Immel, A. Blumberg, J. Merkert, N. Rack, K. Wang, F.
Konstantinidis, J. Truetsch, C. Fernandez, A. Bätz, K. Rösch, M.
Steiner, W. Poh, Y. Shen, R. Wagner, F. Hauser, D. Strutz, J. Villa, G.
Stepanov, H. Caesar, Ö. Ş. Taş, F. Bieder, J.-H. Pauls, and C. Stiller.
The Road Ahead in Autonomous Driving: The KITScenes Multimodal Dataset.
arXiv:2606.02956, 2026.

\end{otherlanguage}

\bibitem{ref-lanelet2}
\begin{otherlanguage}{english}

F. Poggenhans, J.-H. Pauls, J. Janosovits, S. Orf, M. Naumann, F. Kuhnt,
and M. Mayr. Lanelet2: A High-Definition Map Framework for the Future of
Automated Driving. In \emph{Proc. 21st IEEE ITSC}, pp. 1672--1679, 2018.
doi:10.1109/ITSC.2018.8569929.

\end{otherlanguage}

\bibitem{ref-chauffeurnet}
\begin{otherlanguage}{english}

M. Bansal, A. Krizhevsky, and A. Ogale. ChauffeurNet: Learning to Drive
by Imitating the Best and Synthesizing the Worst. In \emph{Proc.
Robotics: Science and Systems (RSS)}, 2019.
doi:10.15607/RSS.2019.XV.031.

\end{otherlanguage}

\bibitem{ref-vectornet}
\begin{otherlanguage}{english}

J. Gao, C. Sun, H. Zhao, Y. Shen, D. Anguelov, C. Li, and C. Schmid.
VectorNet: Encoding HD Maps and Agent Dynamics From Vectorized
Representation. In \emph{Proc. IEEE/CVF CVPR}, pp. 11525--11533, 2020.

\end{otherlanguage}

\bibitem{ref-lanegcn}
\begin{otherlanguage}{english}

M. Liang, B. Yang, R. Hu, Y. Chen, R. Liao, S. Feng, and R. Urtasun.
Learning Lane Graph Representations for Motion Forecasting. In
\emph{Computer Vision -- ECCV 2020}, LNCS 12347, pp. 541--556, Springer,
2020. doi:10.1007/978-3-030-58536-5\_32.

\end{otherlanguage}

\bibitem{ref-admlp}
\begin{otherlanguage}{english}

J.-T. Zhai, Z. Feng, J. Du, Y. Mao, J.-J. Liu, Z. Tan, Y. Zhang, X. Ye,
and J. Wang. Rethinking the Open-Loop Evaluation of End-to-End
Autonomous Driving in nuScenes. arXiv:2305.10430, 2023.

\end{otherlanguage}

\bibitem{ref-egostatus}
\begin{otherlanguage}{english}

Z. Li, Z. Yu, S. Lan, J. Li, J. Kautz, T. Lu, and J. M. Alvarez. Is Ego
Status All You Need for Open-Loop End-to-End Autonomous Driving? In
\emph{Proc. IEEE/CVF CVPR}, pp. 14864--14873, 2024. arXiv:2312.03031.

\end{otherlanguage}

\bibitem{ref-pdm}
\begin{otherlanguage}{english}

D. Dauner, M. Hallgarten, A. Geiger, and K. Chitta. Parting with
Misconceptions about Learning-based Vehicle Motion Planning. In
\emph{Proc. 7th Conference on Robot Learning (CoRL)}, PMLR 229, pp.
1268--1281, 2023.

\end{otherlanguage}

\bibitem{ref-navsim}
\begin{otherlanguage}{english}

D. Dauner, M. Hallgarten, T. Li, X. Weng, Z. Huang, Z. Yang, H. Li, I.
Gilitschenski, B. Ivanovic, M. Pavone, A. Geiger, and K. Chitta. NAVSIM:
Data-Driven Non-Reactive Autonomous Vehicle Simulation and Benchmarking.
In \emph{Advances in Neural Information Processing Systems 37 (NeurIPS
2024), Datasets and Benchmarks Track}, 2024. arXiv:2406.15349.

\end{otherlanguage}

\bibitem{ref-hidden}
\begin{otherlanguage}{english}

B. Jaeger, K. Chitta, and A. Geiger. Hidden Biases of End-to-End Driving
Models. In \emph{Proc. IEEE/CVF ICCV}, pp. 8240--8249, 2023.
arXiv:2306.07957.

\end{otherlanguage}

\bibitem{ref-resnet}
\begin{otherlanguage}{english}

K. He, X. Zhang, S. Ren, and J. Sun. Deep Residual Learning for Image
Recognition. In \emph{Proc. IEEE CVPR}, pp. 770--778, 2016.

\end{otherlanguage}

\end{thebibliography}
\clearpage
\appendix
\subsection{Appendix}\label{appendix}

\subsubsection{Reproducibility}\label{reproducibility}

The evaluated model configuration identifier is
bevformer\_v2\_t8\_split\_navigation\_v5. The camera pathway is
initialized from the official BEVFormer V2 R50 T8 checkpoint (SHA-256
prefix 5585bc4d3ff8; weight license NOASSERTION; trained on nuScenes
under CC-BY-NC-SA-4.0), loaded with the BEV queries and the row/column
positional embeddings resized to 300×200, the camera embeddings matched
to the number of views, and the detection and perspective heads omitted.
The parameter counts reported here were obtained by instantiating this
configuration with 6 views (\hyperref[tab-params]{}). The version
identifiers of the rasterizer, evaluator, and integrator are
navigation\_rasterizer\_v1, reactive\_open\_loop\_metrics\_v1, and
semi\_implicit\_unicycle\_v1. KITScenes originates from
KIT-MRT/KITScenes-Multimodal on Hugging Face (data revision
6fde0034\ldots, SDK revision 7765cdec\ldots). The content-hash prefixes
of the external evaluation index and of the Val and Test manifests are
be1cec5f77a7, 4a12ede73949, and b5dbe72c8d2c. Each internal-validation
evaluation report is recorded with its hash in the model registry, and
we verified that the internal-validation numbers reported here match the
records in the public experiment registry.

\subsubsection{Evaluation identifiers}\label{evaluation-identifiers}

The external evaluations are KITScenes Val v3.5 (role
official\_val\_camera\_map\_route, input track camera\_map\_route, 117
scenes, 11,035 samples) and KITScenes Test v1.0 (role
official\_test\_camera\_only\_missing\_map\_route, input track
camera\_only\_missing\_map\_route, 206 scenes, 23,690 samples). The
inference precision policy is CUDA BF16 automatic mixed precision. The
live Val evaluation policy declares a route-zero pass using a reused
precomputed image BEV, and the publication code requires aggregate
route\_zero\_sample\_count and route\_zero\_trajectory\_delta\_m values.
This indicates that an aggregate route-zero pass was executed upstream.
However, per-sample counterfactual controls and errors are absent from
the external replay artifacts, and the aggregate delta was not included
in the result packet available to this paper. It is therefore neither
reported nor treated as the requested three-condition A/B/C ablation
with per-scene deltas and confidence intervals. Test marks the policy
not applicable. The internal-validation protocols are
reactive\_trajectory\_route\_validation\_6p4s\_v1 (nuPlan) and
kitscenes\_internal\_trajectory\_route\_validation\_6p4s\_v1
(KITScenes).

\subsubsection{Quality gates}\label{quality-gates}

The eval\_gate\_pass field in the registry metadata is false for all
four checkpoints. The repository\textquotesingle s legacy evaluation
workflow contains a gate defined as "6.4 s ADE \textless{} 2.0 m and FDE
\textless{} 4.0 m," but our evidence does not confirm that the registry
field was computed under this definition. We therefore do not speculate
about why the gate was not passed and report only the recorded values.

\subsubsection{Breakdown of comfort
violations}\label{breakdown-of-comfort-violations}

{\def\LTcaptype{none} % do not increment counter
\begingroup\scriptsize
\begin{longtable}[]{@{}llll@{}}
\toprule\noalign{}
Component & Threshold & Epoch 5 & Epoch 7 \\
\midrule\noalign{}
\endhead
\bottomrule\noalign{}
\endlastfoot
Longitudinal acceleration & −4.05 / 2.40 m/s² & 0.0042 & 0.0034 \\
Lateral acceleration (v²κ) & 4.89 m/s² & 0.0147 & 0.0126 \\
Yaw rate (vκ) & 0.95 rad/s & 0.0089 & 0.0079 \\
Yaw acceleration & 1.93 rad/s² & 0.1579 & 0.1401 \\
Longitudinal jerk & 4.13 m/s³ & 0.0000 & 0.0000 \\
Jerk magnitude & 8.37 m/s³ & 0.2356 & 0.2442 \\
Any violation & --- & 0.2458 & 0.2534 \\
\end{longtable}
\endgroup
}

\subsubsection{Selection rule for the qualitative figures and occupancy
statistics}\label{selection-rule-for-the-qualitative-figures-and-occupancy-statistics}

We sorted the 140 shard artifacts of KITScenes Val v3.5 listed by the
public dashboard in lexicographic order by name and, for the first 30
shards, retrieved one sample each at the median index position (the
integer part of half the number of elements). For every sample we
retrieved the map and route rasters, the recorded trajectory, the
navigation metadata, the front-camera tile, and the public control
overlays of the four checkpoints (schema v5, one seed), and rolled them
out with Eq. \hyperref[eq-roll]{}. The 30 samples comprise 22 straight,
4 right-turn, 2 left-turn, and 2 unknown maneuvers; 28 with a valid
route, 30 with a valid map, 26 with a visible goal, and 50 valid future
steps in every case. \hyperref[fig-channels]{}shows the sample with the
largest number of non-empty channels, and the six examples in
\hyperref[fig-qual]{}were chosen by the rule stated in the main text,
with no additional selection by error. Across the sweep, the Epoch 5 5 s
FDE ranged from 0.47 to 14.25 m. These 30 samples were not designed to
be a representative sample of the population.

\subsubsection{Note on raster geometry}\label{note-on-raster-geometry}

The public navigation geometry audited in KITScenes v3 is 256×256 cells
at 1.0 m/px (X: −85.5 to 170.5 m, Y: −128 to 128 m) and covers 99.79\%
of the 6.4 s endpoints. The rasters actually received by the evaluated
checkpoints from KITScenes Val v3.5, in contrast, are 450×300 cells at
0.4 m/px (X: −60 to 120 m, Y: −60 to 60 m), which we confirmed from the
navigation metadata of the public dataset artifacts (geometry\_id =
autoe2e-bev-450x300-0p4m-v1) and from the array shapes (14×450×300 and
2×450×300). The navigation encoder bilinearly resamples these onto a
300×200 latent grid, so their extent matches that of the camera BEV. BEV
segmentation evaluation nearest-neighbor-resamples the 1.0 m geometry
labels onto the 0.4 m grid, but we do not report BEV segmentation
numbers in this paper.

\subsubsection{Claim audit}\label{claim-audit}

{\def\LTcaptype{none} % do not increment counter
\begingroup\scriptsize
\begin{longtable}[]{@{}lll@{}}
\toprule\noalign{}
Claim & Type & Evidence \\
\midrule\noalign{}
\endhead
\bottomrule\noalign{}
\endlastfoot
The camera pathway is based on the official BEVFormer V2 R50 T8 and is
frozen during the trajectory-training stage & Implementation, registry &
Initialization code and freeze contract; training configuration
freeze\_bevformer = true \\
BEV grid 300×200, X −60 to 120 m, Y −60 to 60 m, 0.6 m cells &
Implementation & Fixed model configuration; confirmed by
instantiation \\
7 history frames plus the current frame, 0.5 s spacing, 3.5 s; history
at 512 px only & Implementation & T8 contract constants and
history-encoding code \\
The 1024 px front branch is added for the current frame only; the 512 px
front view remains in the all-view encoder & Implementation & Front
residual-branch code; zero-init gate \\
Map 14 ch, route 2 ch, corridor 3.5 m, goal radius 2.0 m &
Implementation & Channel enumeration and rasterizer; array shapes from
real data \\
The evaluated map and route rasters are 450×300 at 0.4 m/px & Registry,
implementation & geometry\_id and array shapes of Val v3.5 samples;
data-contract validation code \\
Deformable fusion uses 8 points, 8 heads, 256 ch, and is identity at
initialization & Implementation & zero-init of the fusion module \\
GRU planner, 16 sample points, 64 steps, acceleration and curvature
outputs & Implementation & Planner configuration num\_points = 16 \\
79,906,522 parameters in total, 1,448,582 trainable & Implementation &
Parameter count from our instantiation; consistent with the repository
record of 79,907,034 (8 views) \\
enable\_world\_model = false, temporal\_memory\_mode = no\_memory;
h\_vis is zero & Implementation, registry & Fixed configuration;
KITScenes sample-generation code \\
Training configuration of the four checkpoints (8 workers, seed 149,
weights 1.0/1.0/0.0, LR) & Registry & Checkpoint registry metadata;
learning-rate records in the public experiment registry \\
eval\_gate\_pass = false (all four checkpoints) & Registry & Registry
metadata; gate definition unconfirmed \\
On Val, Epoch 7 is best at 1 s, while Epoch 5 is best at 2/3/5 s and on
lateral and longitudinal error & Measured &
\hyperref[tab-val-ade]{},\hyperref[tab-latlon]{} \\
On Test, Epoch 5 is best at 1 s and 5 s; on nuPlan, Epoch 5 is best at 2
s and 3 s & Measured & \hyperref[tab-test-ade]{} \\
Zero non-finite predictions, 78.125\% coverage & Measured &
\hyperref[tab-integrity]{} \\
Epoch 7 improves the route-related proxies while degrading 6.4 s error,
drivable-area success rate, and comfort rate & Measured &
\hyperref[tab-internal-kit]{}; matches the values in the public
experiment registry \\
Additional training improves near-horizon accuracy and degrades mid- and
long-horizon accuracy & Interpretation & Rank inversion between 1 s and
2--5 s on Val \\
Checkpoint selection should be multi-objective & Interpretation &
Cross-metric trade-off between Epoch 5 and Epoch 7 \\
Route-reconstruction IoU verifies information preservation, not route
adherence & Interpretation & The head takes the post-gate route
contribution as input; divergence from the adherence rate \\
The Test results do not show that the map is unnecessary &
Interpretation & Difference between the Val and Test scene
populations \\
HD map and route should resolve geometric ambiguity and branching &
Hypothesis & Controlled experiment not run \\
Benefits of separate encoding and residual gated fusion & Hypothesis &
Controlled experiment not run \\
Map and route causally improved planning accuracy & Not claimed & No
same-scene control; counterfactuals not stored; confidence intervals not
computed \\
\end{longtable}
\endgroup
}

\subsubsection{Additional experiments required for causal
claims}\label{additional-experiments-required-for-causal-claims}

\begin{figure}
\centering
\pandocbounded{\includegraphics[keepaspectratio,width=\linewidth]{figures_pdf/fig15_matched_ablation.pdf}}
\caption{Design of the same-scene controlled experiment required to
establish the effect of map and route. On the same 11,035 KITScenes Val
samples and the same frozen checkpoint, the same precomputed camera BEV
features are reused while only the navigation input is switched among A
(camera only), B (camera + HD map), and C (camera + HD map + route). We
do not have results for this experiment.}\label{fig-ablation}
\end{figure}

The same 11,035 KITScenes Val samples and the same frozen checkpoint
would be evaluated in matched fashion under three conditions: A, camera
only; B, camera + HD map; C, camera + HD map + route. The same
precomputed camera BEV features are reused across the three conditions,
and only the navigation input is changed. Switching uses the explicit
map\_valid, route\_valid, and enable\_route\_conditioning gates;
stochastic augmentation is not re-run, and the planner configuration is
held fixed. The reported quantities are ADE/FDE at 1, 2, 3, and 5 s;
mean absolute lateral and longitudinal error; route-corridor adherence
and success rates; the projected-endpoint arc-length proxy;
drivable-area adherence and success rates; the comfort rate and each
violation component; the failure rate and non-finite predictions; and,
where possible, inference latency and memory. Scene-level paired
bootstrap confidence intervals and the distribution of per-scene
differences would be reported. Additional controls would include a
shuffled route taken from a different scene, removal of the goal marker,
removal of the route corridor, map-only corruption and dropout,
off-route and ambiguous-intersection subsets, and stratification by
valid versus invalid map. These results are not included in this paper,
and we do not predict them.

\end{document}
